\documentclass[12pt, twoside]{article}
\input{../preamble}

\fancyhead[LO]{La Scuola d'Italia / Dr.\ Huson / IB Math \\* 22 April 2026}
\fancyhead[RO]{\fbox{\begin{minipage}{6cm}\footnotesize First \& Last name:\\[0.5cm]\end{minipage}}}
\fancyfoot[C]{\thepage}

\begin{document}
\subsubsection*{5.15 Classwork: Logarithms (no calculator)}

\begin{enumerate}[itemsep=0.15cm]

%--- Problem 1: Logarithm evaluation and variable expressions [12 marks] ---
%--- Source: original, modeled on AA SL 2022 Nov P1 Q8(a) and Math SL 2018 Nov P1 Q4 ---
\item[1.] [Maximum mark: 12]

\begin{enumerate}[itemsep=0.2cm]
  \item Calculate the value of each of the following logarithms:
  \begin{enumerate}[itemsep=0.1cm, label=(\roman*)]
    \item $\log_4 \dfrac{1}{64}$
    \item $\log_{25} 5$
    \item $\log_2 \sqrt{32}$ \hfill [6]
  \end{enumerate}

  \item Let $c = \log_3 m$, where $m > 0$. Write down each of the
  following expressions in terms of $c$.
  \begin{enumerate}[itemsep=0.1cm, label=(\roman*)]
    \item $\log_3 m^2$
    \item $\log_3 27m$
    \item $\log_9 m$ \hfill [6]
  \end{enumerate}
\end{enumerate}

\begin{tikzpicture}
  \draw (-0.6,0) rectangle (11,9);
  \foreach \y in {1,2,...,8} \draw[dotted] (0,\y)--(10.5,\y);
\end{tikzpicture}

\newpage

%--- Problem 2: Log equation with base conversion [7 marks] ---
%--- Source: original, modeled on Math SL 2019 May P1 TZ2 Q6 ---
\item[2.] [Maximum mark: 7]

Solve $\log_2(x - 3) = \log_4(x + 9)$.

\begin{tikzpicture}
  \draw (-0.6,0) rectangle (11,14);
  \foreach \y in {1,2,...,13} \draw[dotted] (0,\y)--(10.5,\y);
\end{tikzpicture}

\newpage

\noindent Answer all questions on lined paper.
Please start each question on a new page.

\medskip

%--- Problem 3: Log series geometric/arithmetic [15 marks] ---
%--- Source: AA SL 2022 May P1 TZ1 Q08 ---
\item[3.] [Maximum mark: 15]

Consider the series $\ln x + p \ln x + \dfrac{1}{3} \ln x + \ldots$,
where $x \in \mathbb{R}$, $x > 1$ and $p \in \mathbb{R}$, $p \neq 0$.

\begin{enumerate}[itemsep=0.2cm]
  \item Consider the case where the series is geometric.
  \begin{enumerate}[itemsep=0.1cm, label=(\roman*)]
    \item Show that $p = \pm \dfrac{1}{\sqrt{3}}$.
    \item Given that $p > 0$ and $S_{\infty} = 3 + \sqrt{3}$,
    find the value of $x$. \hfill [5]
  \end{enumerate}

  \item Now consider the case where the series is arithmetic
  with common difference $d$.
  \begin{enumerate}[itemsep=0.1cm, label=(\roman*)]
    \item Show that $p = \dfrac{2}{3}$.
    \item Write down $d$ in the form $k \ln x$, where
    $k \in \mathbb{R}$.
    \item The sum of the first $n$ terms of the series is
    $-3 \ln x$. Find the value of $n$. \hfill [10]
  \end{enumerate}
\end{enumerate}

\bigskip

%--- Problem 4: Exponential function and inverse [14 marks] ---
%--- Source: Math SL 2019 Nov P1 Q10 ---
\item[4.] [Maximum mark: 14]

Let $g(x) = p^x + q$, for $x, \, p, \, q \in \mathbb{R}$, $p > 1$.
The point $A(0, a)$ lies on the graph of $g$.

Let $f(x) = g^{-1}(x)$. The point $B$ lies on the graph
of $f$ and is the reflection of point $A$ in the line $y = x$.

\begin{enumerate}[itemsep=0.15cm]
  \item Write down the coordinates of $B$. \hfill [2]
\end{enumerate}

The line $L_1$ is tangent to the graph of $f$ at $B$.

\begin{enumerate}[itemsep=0.15cm]
  \item[(b)] Given that $f'(a) = \dfrac{1}{\ln p}$, find the equation
  of $L_1$ in terms of $x$, $p$ and $q$. \hfill [5]
\end{enumerate}

The line $L_2$ is tangent to the graph of $g$ at $A$ and has
equation $y = (\ln p)\, x + q + 1$.

The line $L_2$ passes through the point $(-2, -2)$.

The gradient of the normal to $g$ at $A$ is
$\dfrac{1}{\ln \frac{1}{3}}$.

\begin{enumerate}[itemsep=0.15cm]
  \item[(c)] Find the equation of $L_1$ in terms of $x$. \hfill [7]
\end{enumerate}

\newpage

%--- Problem 5: Kinematics, displacement function [16 marks] ---
%--- Source: AA SL 2022 Nov P2 Q7 ---
\item[5.] [Maximum mark: 16]

A particle moves in a straight line. Its displacement, $s$ metres,
from a fixed point $P$ at time $t$ seconds is given by
$s(t) = 3(t + 2)^{\cos t}$, for $0 \leq t \leq 6.8$, as shown in the
following graph.

\begin{center}
\begin{tikzpicture}[x=1.2cm, y=0.18cm]
  \draw[->] (-0.3, 0) -- (7.3, 0) node[right] {$t$};
  \draw[->] (0, -1) -- (0, 28) node[above] {$s$};
  \node[below left] at (0, 0) {$0$};
  \draw (6.8, -0.8) -- (6.8, 0.8);
  \node[below] at (6.8, -0.8) {$6.8$};
  \draw[thick, samples=200, domain=0:6.8, smooth] 
    plot (\x, {3*(\x+2)^(cos(deg(\x)))});
\end{tikzpicture}
\end{center}

\begin{enumerate}[itemsep=0.15cm]
  \item Find the particle's initial displacement from the point $P$.
  \hfill [2]
  \item[(b)] Find the particle's velocity when $t = 2$. \hfill [2]
  \item[(c)] Determine the intervals of time when the particle is
  moving away from  \\[0.25cm]
  the point $P$. \hfill [5]
\end{enumerate}

The acceleration of the particle is zero when $t = b$ and $t = c$,
where $b < c$.

\begin{enumerate}[itemsep=0.15cm]
  \item[(d)] Find the value of $b$ and the value of $c$. \hfill [4]
  \item[(e)] Find the total distance travelled by the particle for
  $b \leq t \leq c$. \hfill [3]
\end{enumerate}

\end{enumerate}
\end{document}
